% =====================================================================
% CHAPTER 4: DATASET DESCRIPTION / DETAILS
% =====================================================================
\chapter{Dataset Description / Details}
\label{ch:dataset}

\section{Overview}
\label{sec:dataset-overview}

\textbf{No formal benchmark dataset is used by the current implementation of TrackVision.} The system is designed for real-time inference on live camera streams or user-provided video input, not for offline evaluation on standardized MOT benchmarks (MOT17 \cite{milan2016mot16}, MOT20 \cite{dendorfer2020mot20}, KITTI \cite{geiger2012kitti}, TAO \cite{dave2020tao}, DanceTrack \cite{sun2022dancetrack}).

This design choice reflects the project's primary goal: \textbf{a deployable, client-side MOT application} rather than a benchmark submission. Consequently, standardized quantitative metrics (MOTA, IDF1, HOTA, ID switches) cannot be computed and are not reported in this work.

\section{Input Sources}
\label{sec:input-sources}

TrackVision accepts two input modalities:

\begin{enumerate}
    \item \textbf{Live Camera Stream:} Captured via \texttt{navigator.mediaDevices.getUserMedia()} with constraints \texttt{\{ video: \{ facingMode: 'environment', width: \{ ideal: 1280 \}, height: \{ ideal: 720 \} \}, audio: false \}}. The rear-facing camera is preferred for mobile scenarios.
    \item \textbf{User-Provided Video Files:} The application supports drag-and-drop or file selection of browser-compatible video formats (MP4, WebM, MOV, etc.). The \texttt{<video>} element loads the file via \texttt{URL.createObjectURL()}.
\end{enumerate}

Both sources feed into the same processing pipeline (Figure~\ref{fig:pipeline}).

\section{Video Characteristics}
\label{sec:video-characteristics}

Since the input is user-dependent, the following parameters are \textit{not fixed} but are handled dynamically by the pipeline:

\begin{table}[htbp]
\centering
\caption{Input video parameter handling}
\label{tab:video-params}
\begin{tabular}{@{}ll@{}}
\toprule
\textbf{Parameter} & \textbf{Handling} \\
\midrule
Resolution & Variable (camera: up to 1280×720; file: any); letterboxed to 640×640 internally \\
Frame Rate & Variable (typically 30 FPS); tracker assumes $\Delta t = 1/30$ s \\
Color Space & RGB (browser canvas); converted to CHW float32 for inference \\
Aspect Ratio & Preserved via letterboxing (padding); scale/offset metadata retained \\
Duration & Unlimited (5-minute rolling history window for storage) \\
\bottomrule
\end{tabular}
\end{table}

\section{Preprocessing Pipeline}
\label{sec:preprocessing}

Every frame undergoes the following deterministic preprocessing before inference:

\begin{enumerate}
    \item \textbf{Capture:} \texttt{OffscreenCanvas.drawImage(video, 0, 0, 640, 640)} with letterboxing (aspect-preserving fit with padding).
    \item \textbf{Metadata:} The capture returns \texttt{scale} (min(640/w, 640/h)), \texttt{offsetX}, \texttt{offsetY} (padding offsets).
    \item \textbf{Planar Conversion:} RGBA $\to$ float32 CHW (3 × 640 × 640), normalized by 1/255.
    \item \textbf{Buffer Transfer:} \texttt{ArrayBuffer} transferred to Web Workers (detached, zero-copy).
\end{enumerate}

This preprocessing is implemented in \texttt{src/hooks/useOffscreenCanvas.ts} and the detection workers.

\section{Object Classes}
\label{sec:object-classes}

\begin{itemize}
    \item \textbf{Fast Mode (YOLOv8n):} 80 fixed COCO classes (\texttt{COCO\_CLASSES} in \texttt{src/workers/workerUtils.ts}). Examples: person, bicycle, car, motorbike, bus, truck, traffic light, dog, cat, laptop, phone, bottle, cup, chair, book, clock, etc.
    \item \textbf{Open Mode (YOLO-World):} Dynamic, user-defined concepts. The user enters comma-separated text prompts (e.g., ``laptop, phone, headphones, water bottle''). Each concept is formatted as \texttt{``a photo of a \{concept\}"}, tokenized by the in-browser CLIP BPE tokenizer, and encoded by the CLIP text encoder to produce \texttt{text\_features [1, C, 512]}.
\end{itemize}

\section{Annotations}
\label{sec:annotations}

\textbf{No ground-truth annotations are available.} The system operates in a purely \textit{inference-only} mode. There is no training, validation, or test split. This is a fundamental limitation for quantitative MOT evaluation.

\section{Test Videos Used During Development}
\label{sec:test-videos}

During manual development and qualitative testing, the following video sources were used:

\begin{itemize}
    \item Live webcam (laptop integrated camera, USB webcam).
    \item Mobile phone camera (via \texttt{dev:https} LAN access).
    \item Sample video: \texttt{one-by-one-person-detection.mp4} (Intel IoT DevKit sample, $\sim$3.3 MB, 1920×1080, 30 FPS, single person walking).
\end{itemize}

These are \textbf{not benchmark datasets} and were used solely for functional verification and UI testing.

\section{Summary Table}
\label{sec:dataset-summary}

\begin{table}[htbp]
\centering
\caption{Dataset description summary}
\label{tab:dataset-summary}
\begin{tabular}{@{}ll@{}}
\toprule
\textbf{Parameter} & \textbf{Value} \\
\midrule
Dataset Name & \textit{None (live camera / user video only)} \\
Source & Browser \texttt{getUserMedia()} / File input \\
Video Count & Unlimited (user-dependent) \\
Resolution & Variable; processed at 640×640 letterboxed \\
FPS & Variable; tracker $\Delta t = 1/30$ s \\
Object Classes & COCO-80 (fast) / User-defined (open) \\
Annotation Format & None (inference only) \\
Train/Test Split & N/A \\
Ground Truth & Not available \\
\bottomrule
\end{tabular}
\end{table}

\section{Implication for Evaluation}
\label{sec:eval-implication}

The absence of a benchmark dataset with ground-truth trajectories means that \textbf{standardized MOT metrics cannot be computed}. Chapter~\ref{ch:results} discusses the evaluation methodology actually used (unit tests, microbenchmarks, qualitative observation) and proposes a future work path for standardized evaluation.